% ===== PRÉAMBULE CANONIQUE — à coller VERBATIM en tête de CHAQUE .tex =====
\documentclass[11pt,a4paper]{article}
\usepackage[utf8]{inputenc}\usepackage[T1]{fontenc}\usepackage{lmodern}
\usepackage[french]{babel}
\usepackage{amsmath,amssymb,mathtools}\usepackage{icomma}
\usepackage{siunitx}\sisetup{locale=FR}  % \qty{}{}, \si{}, \ang{}
\usepackage{xcolor}\usepackage{colortbl}
\usepackage{tikz}\usepackage{pgfplots}\pgfplotsset{compat=1.18}
\usepackage[siunitx,europeanresistors]{circuitikz} % circuits (élec.)
\usepackage[most]{tcolorbox}\usepackage{enumitem}\usepackage{array,booktabs}
\usepackage[a4paper,margin=2.2cm]{geometry}
\usetikzlibrary{arrows.meta,calc,angles,quotes,positioning,patterns,%
  backgrounds,shapes.geometric,decorations.pathmorphing,decorations.markings,intersections}
\definecolor{primary}{RGB}{222,49,99}\definecolor{primarydark}{RGB}{150,22,55}
\definecolor{defcol}{RGB}{0,114,178}\definecolor{methcol}{RGB}{52,92,148}
\definecolor{excol}{RGB}{0,135,90}\definecolor{attncol}{RGB}{200,40,55}
\tcbset{boxbase/.style={enhanced,breakable,boxrule=0.9pt,arc=2.5pt,
  left=3mm,right=3mm,top=2.3mm,bottom=2.3mm,fonttitle=\bfseries,coltitle=white}}
% --- boîtes TUTORIEL (noms EXACTS, ne pas renommer) ---
\newtcolorbox{cle}[1][]{boxbase,colback=defcol!6,colframe=defcol,
  title={\textsc{À retenir}\ifx\relax#1\relax\relax\else\textmd{ : \itshape #1}\fi}}
\newtcolorbox{methode}[1][]{boxbase,colback=methcol!7,colframe=methcol,sharp corners,
  title={\textsc{Méthode}\ifx\relax#1\relax\relax\else\textmd{ --- \itshape #1}\fi}}
\newtcolorbox{exemplebox}[1][]{boxbase,colback=excol!5,colframe=excol,
  title={\textsc{Exemple}\ifx\relax#1\relax\relax\else\textmd{ : \itshape #1}\fi}}
\newtcolorbox{piege}[1][]{boxbase,colback=attncol!6,colframe=attncol,
  title={\textsc{Piège}\ifx\relax#1\relax\relax\else\textmd{ : \itshape #1}\fi}}
\newtcolorbox{remarque}[1][]{boxbase,colback=black!4,colframe=black!45,
  title={\textsc{Remarque}\ifx\relax#1\relax\relax\else\textmd{ : \itshape #1}\fi}}
% --- boîtes EXERCICE / CORRIGÉ (noms EXACTS, auto-comptées) ---
\newtcolorbox[auto counter]{exo}[1][]{boxbase,colback=primary!4,colframe=primary,
  title={\textsc{Exercice}~\thetcbcounter\ifx\relax#1\relax\relax\else\textmd{ --- \itshape #1}\fi}}
\newtcolorbox[auto counter]{corrige}[1][]{boxbase,colback=excol!4,colframe=excol,
  title={\textsc{Corrigé}~\thetcbcounter\ifx\relax#1\relax\relax\else\textmd{ --- \itshape #1}\fi}}
\newcommand{\vv}[1]{\vec{#1}}\newcommand{\ds}{\displaystyle}
\newcommand{\R}{\mathbb{R}}\newcommand{\N}{\mathbb{N}}\newcommand{\Z}{\mathbb{Z}}
% (un \usepackage{fancyhdr}+en-tête est possible ; il est ignoré côté web)
% =========================================================================
\begin{document}
\sisetup{per-mode=symbol}

\begin{exo}[deux charges de même signe]
Le spectre ci-dessous représente deux charges ponctuelles $A$ et $B$. Les lignes s'écartent l'une de
l'autre et aucune ne relie directement les deux charges.
\begin{center}
\begin{tikzpicture}[scale=0.75,>=Stealth,font=\scriptsize,
    ch/.style={circle,fill=attncol,minimum size=6.5mm,inner sep=0pt,text=white,font=\bfseries\scriptsize}]
  \node[ch] (A) at (-1.8,0) {$A$}; \node[ch] (B) at (1.8,0) {$B$};
  \foreach \a in {150,180,210,120,240}{\draw[->,excol,line width=0.7pt] (-1.8,0)--++(\a:1.7);}
  \foreach \a in {30,0,-30,60,-60}{\draw[->,excol,line width=0.7pt] (1.8,0)--++(\a:1.7);}
  \draw[->,excol,line width=0.7pt] (-1.8,0) .. controls (-0.9,1.3) and (0.9,1.3) .. (1.8,0);
  \draw[->,excol,line width=0.7pt] (-1.8,0) .. controls (-0.9,-1.3) and (0.9,-1.3) .. (1.8,0);
  \node[circle,draw=primary,fill=white,inner sep=1.4pt] at (0,0) {}; \node[primary,font=\scriptsize,above] at (0,0.12) {$N$};
\end{tikzpicture}
\end{center}
\begin{enumerate}[label=\alph*)]
  \item Les charges $A$ et $B$ sont-elles de même signe ou de signes opposés ? Justifie.
  \item Si $A$ et $B$ sont égales, où se situe le \emph{point neutre} $N$ (où $\vv{E}=\vv{0}$) ?
\end{enumerate}
\end{exo}

\begin{exo}[comparer deux charges par la densité]
Deux charges positives $A$ et $B$ émettent respectivement \textbf{8} et \textbf{4} lignes de champ.
\begin{center}
\begin{tikzpicture}[scale=0.7,>=Stealth,font=\scriptsize,
    ch/.style={circle,fill=attncol,minimum size=6mm,inner sep=0pt,text=white,font=\bfseries\scriptsize}]
  \begin{scope}[xshift=-3cm]
    \foreach \a in {0,45,...,315}{\draw[->,excol,line width=0.7pt] (0,0)--(\a:1.5);}
    \node[ch] at (0,0) {$A$}; \node[primarydark] at (0,-2.1) {8 lignes};
  \end{scope}
  \begin{scope}[xshift=3cm]
    \foreach \a in {0,90,180,270}{\draw[->,excol,line width=0.7pt] (0,0)--(\a:1.5);}
    \node[ch] at (0,0) {$B$}; \node[primarydark] at (0,-2.1) {4 lignes};
  \end{scope}
\end{tikzpicture}
\end{center}
\begin{enumerate}[label=\alph*)]
  \item Quelle charge est la plus grande ?
  \item Le nombre de lignes étant proportionnel à la charge, quel est le rapport $q_A/q_B$ ?
\end{enumerate}
\end{exo}

\begin{exo}[champ dans un condensateur]
Un condensateur plan a ses armatures séparées de $d = \SI{5}{\centi\metre}$ et soumises à une tension
$U_{AB} = \SI{200}{\volt}$ ($V_A > V_B$).
\begin{enumerate}[label=\alph*)]
  \item Calcule l'intensité du champ uniforme entre les armatures.
  \item Vers quelle armature ce champ est-il orienté ?
\end{enumerate}
\end{exo}

\begin{exo}[une charge lâchée dans le champ]
Entre les armatures d'un condensateur, le champ $\vv{E}$ est uniforme et dirigé \textbf{vers le bas}.
On y abandonne, sans vitesse initiale, une petite charge (on néglige la pesanteur).
\begin{enumerate}[label=\alph*)]
  \item Dans quel sens part une charge \textbf{positive} ?
  \item Dans quel sens part une charge \textbf{négative} ?
\end{enumerate}
\end{exo}

\begin{exo}[uniforme ou radial ?]
Voici deux spectres.
\begin{center}
\begin{tikzpicture}[scale=0.7,>=Stealth,font=\scriptsize,
    ch/.style={circle,fill=attncol,minimum size=6mm,inner sep=0pt,text=white,font=\bfseries\scriptsize}]
  \begin{scope}[xshift=-3.4cm]
    \foreach \x in {-1.4,-0.7,0,0.7,1.4}{\draw[->,excol,line width=0.8pt] (\x,1.3)--(\x,-1.3);}
    \node[primarydark] at (0,-1.9) {spectre I};
  \end{scope}
  \begin{scope}[xshift=3.4cm]
    \foreach \a in {0,45,...,315}{\draw[->,excol,line width=0.8pt] (0,0)--(\a:1.5);}
    \node[ch] at (0,0) {$+$};
    \node[primarydark] at (0,-1.9) {spectre II};
  \end{scope}
\end{tikzpicture}
\end{center}
\begin{enumerate}[label=\alph*)]
  \item Lequel correspond à un champ \textbf{uniforme} (condensateur) ? Justifie.
  \item Dans ce spectre uniforme, comment se compare l'intensité $E$ en deux points différents ?
\end{enumerate}
\end{exo}

\end{document}
